% !TEX root = ../main.tex
%----------------------------------------------------------
% 02-introduction.tex — Section 1: Introduction
%----------------------------------------------------------

\section{Introduction}\label{sec:introduction}

Consider two forecasts, each assigning equal odds.  A statistician
estimates a fair coin at 50:50 after one hundred flips---confident,
grounded in data.  A pundit predicts 50:50 for a sports match played in
extreme weather between teams that have never met---uncertain, grounded
in ignorance.  Probability theory assigns both the same distribution;
possibility theory does not.  The coin forecast is \textit{normal}
($\Pi_{\max} = 1$, $\ign = 0$), fully endorsing its estimate; the match
forecast is \textit{subnormal} ($\Pi_{\max} \ll 1$, high $\ign$),
encoding that equal odds arose from absence of evidence rather than
weight of evidence.

Hazard forecasting increasingly relies on probabilistic methods to
communicate uncertainty, and for good reason: calibrated probability
distributions support principled decision-making and admit rigorous
verification through proper scoring rules.  Yet probability demands
precise quantification---each outcome receives a single number
representing relative frequency or degree of belief, and those numbers
must sum to unity.  Not all forecasting systems produce output that
satisfies these constraints.  The Storm Prediction Center (SPC)
convective outlooks, for instance, assign ordered categorical labels
(\spc{NONE} through \spc{HIGH}) that express graded plausibility rather
than calibrated frequency.  SPC issues crisp categorical outlooks drawn
on a map; this paper proposes a hypothetical possibilistic alternative
using these familiar categories as the running example, treating the
outlook at a single point (Norman, Oklahoma) as the observation.
Possibility theory \citep{Dubois1988-nh}
provides a mathematical framework for exactly this kind of statement:
graded compatibility of outcomes with available evidence, without the
additivity requirement that probability imposes.  Despite a mature
theoretical literature, no verification methodology exists for
possibilistic forecasts.  This paper fills that gap.

% --- Why possibility theory matters for hazard verification ---

A probabilistic forecast captures one dimension of uncertainty: given
sufficient evidence, each outcome receives a precise weight.  A
possibility distribution captures something richer.  For any event~$A$,
the interval $[\nec(A),\, \poss(A)]$ brackets the range of
probabilities consistent with the available evidence
\citep{Dubois2006-MISSING}.  When evidence is abundant, the interval
collapses and possibility reduces to a point probability; when evidence
is scarce, the interval widens, faithfully representing the forecaster's
epistemic state.  This bracketing property makes possibility theory a
strict generalization of probability rather than a competing paradigm.

Possibility distributions arise naturally in several operational
settings.  Expert-rule systems and fuzzy inference engines produce
membership functions that are possibility distributions by construction
\citep{Klir1995-va}.  Imprecise-probability frameworks generate upper
and lower bounds on event probabilities that map directly onto the
possibility--necessity duality.  In severe-weather forecasting, a
forecaster may wish to express that \spc{MDT} risk is plausible but not
certain, while acknowledging that the mesoscale environment is poorly
sampled---a statement that maps cleanly onto a subnormal possibility
distribution but awkwardly onto a probability.  Assigning
$\pi(\spc{MDT}) = 0.7$ says ``\spc{MDT} is fairly compatible with
available evidence''; it does not say ``the long-run frequency of
\spc{MDT} events is 70\%.''

The distinction becomes sharpest for subnormal distributions, where the
maximum possibility value falls below unity: $\max_\omega \pi(\omega) <
1$.  A subnormal distribution encodes a signal that probabilistic
forecasts cannot represent: the system does not fully endorse any
outcome.  This possibilistic ignorance, denoted $\ign = 1 - \max(\pi)$,
quantifies the gap between the system's best-supported outcome and full
endorsement.  Probabilistic frameworks have no direct counterpart---an
ensemble that assigns low probability to every member still produces a
distribution that sums to one, absorbing the uncertainty into the
shape of the distribution rather than flagging it as a distinct signal.

% --- The "ignorance erasure" problem ---

Classical possibility theory requires normalization: at least one
outcome must receive $\pi(\omega) = 1$ for the necessity measure to
be well-defined \citep{Dubois1988-nh}.  In a verification context,
normalizing a subnormal distribution---rescaling so that $\max(\pi) =
1$---destroys precisely the ignorance signal that distinguishes a
confident forecast from a tentative one.  A system that produces
$\pi = (0.05,\, 0.2,\, 0.4,\, 0.6,\, 0.1,\, 0.0)$, indicating 40\% ignorance,
becomes indistinguishable after normalization from a system that
confidently assigns $\pi = (0.08,\, 0.33,\, 0.67,\, 1.0,\, 0.17,\, 0.0)$.
Both carry the same shape, but the first system was honestly reporting
that it lacked evidence.

This pattern is analogous to a familiar problem in ensemble
verification.  When all ensemble members cluster in a low-probability
region, post-processing must still produce calibrated probabilities that
sum to unity; the ensemble's collective uncertainty is absorbed into the
distribution rather than flagged.  We call this \textit{ignorance
erasure}: the systematic destruction of subnormality information through
normalization or probability conversion.  Ignorance erasure is, we
argue, the central obstacle to adopting possibilistic methods in
operational verification, because it conflates ``the system thinks
\spc{ENH} is most likely'' with ``the system is confident.''

% --- Brief literature review ---

Possibility theory originates with \citet{Zadeh1978-je}, who interpreted fuzzy
membership functions as possibility distributions.  It was formalized
into a full uncertainty calculus by \citet{Dubois1988-nh}.
\citet{Dubois2006-MISSING} surveys the mature theory, emphasizing the
bracketing property $\nec(A) \leq P(A) \leq \poss(A)$ and its
applications to risk analysis under incomplete information.
\citet{Klir1995-va} situates possibility within a broader taxonomy of
generalized information theories, clarifying the relationship between
possibilistic and probabilistic uncertainty measures.

A parallel literature addresses the verification of probabilistic
forecasts through proper scoring rules.  \citet{Roulston2002-eq}
introduced the information-gain framework, measuring forecast value as
the reduction in Shannon entropy relative to a climatological baseline.
\citet{Lawson2024-bu} formalized the information-gain decomposition for
atmospheric forecasts, separating overall skill into uncertainty,
discrimination, and reliability components.  These tools are well suited
to probability vectors but cannot be applied directly to possibility
distributions, which violate the additivity axiom.

The Dempster--Shafer theory of evidence \citep{Shafer1976-MISSING}
provides an adjacent framework in which belief and plausibility functions
generalize probability.  \citet{Smets1990-MISSING} introduced the
pignistic transformation---from the Latin \textit{pignus}, a wager---to
convert belief functions into decision-ready probabilities.  The theory
of imprecise probabilities \citep{Walley1991-MISSING} further develops
the idea of interval-valued uncertainty representations.  All three frameworks share
the intuition that a single probability may not adequately represent an
agent's epistemic state, but none provides a verification methodology
tailored to subnormal possibility distributions.

Despite the growing use of possibility-like representations in
expert systems, fuzzy inference, and imprecise-probability applications,
no systematic verification framework addresses subnormal possibility
distributions.  Existing approaches either normalize the distribution---
erasing the ignorance signal---or treat possibility values as
probabilities, violating the max-additivity axiom that distinguishes
possibility from probability.  This paper develops a unified framework
with three components: a five-number scorecard for native possibilistic
evaluation, a possibility-to-probability bridge that preserves ignorance
as an explicit outcome, and an interval log score derived from the
possibility--necessity duality.

% --- Paper outline ---

The remainder of this paper is organized as follows.
Section~\ref{sec:poss_primer} introduces the mathematical foundations of
possibility theory, emphasizing subnormal distributions and the
three-component uncertainty framework.
Section~\ref{sec:pignistic_bridge} develops the possibility-to-probability
bridge that converts possibilities to probabilities without erasing
ignorance, and reviews the information-gain decomposition.
Section~\ref{sec:native_verification} formalizes native possibilistic
verification: the five-number scorecard, interval log score, imprecise
ranked logarithmic score, and possibilistic climatology.
Section~\ref{sec:tripartite_value} establishes three hypothesis tests
demonstrating that each uncertainty component (possibility, ignorance,
conditional necessity) independently contributes verification value.
Section~\ref{sec:worked_examples} walks through three worked examples
using SPC convective outlook categories (\spc{NONE} through \spc{HIGH}).
Section~\ref{sec:discussion} discusses connections to Dempster--Shafer
theory, applicability beyond severe weather, and future directions.
